\chapter{La substitution nucléophile}\label{ch:b1:nucleophilic-substitution}

Prenons un flacon de \cip{S}-2-bromobutane, un énantiomère unique qui
fait tourner la lumière polarisée. Chauffé avec de l'hydroxyde de sodium
dans un \omterm{def:b1:intermolecular-forces-solvents:solvent-classes}{solvant polaire} aprotique, il donne du butan-2-ol qui est lui
aussi un énantiomère unique --- mais l'autre, \cip{R}~: la réaction a
retourné la molécule comme un parapluie dans le vent. Laissé au contraire
dans de l'eau tiède, le même bromure donne du butan-2-ol en partie
racémique. Une même transformation globale, un groupe OH remplaçant un
atome Br, suit deux chemins différents. Ce chapitre établit les deux
mécanismes à partir de leurs \omterm{def:b1:rate-laws:order}{lois de vitesse}, explique leur
stéréochimie et donne les règles qui permettent de choisir entre eux ---
c'est la première utilisation complète des outils des chapitres de
cinétique et d'effets électroniques.

\begin{recall}
Les \omterm{def:b1:rate-laws:order}{lois de vitesse}, les \omterm{def:b1:elementary-steps:elementary-step}{actes élémentaires}, l'\omterm{def:b1:elementary-steps:rds}{étape cinétiquement
déterminante} et l'\omterm{def:b1:elementary-steps:qssa}{approximation des états quasi stationnaires} sont aux
\cref{ch:b1:rate-laws,ch:b1:elementary-steps}~; les \omterm{def:b1:stereochemistry-in-depth:isomers}{configurations}, les
\omterm{def:b1:stereochemistry-in-depth:representations}{représentations de Cram} et l'\omterm{def:b1:stereochemistry-in-depth:optical-activity}{activité optique} au
\cref{ch:b1:stereochemistry-in-depth}~; les \omterm{def:b1:electronic-effects:nucleophile}{nucléophiles}, les
\omterm{def:b1:electronic-effects:nucleophile}{électrophiles}, les \omterm{def:b1:electronic-effects:nucleophile}{groupes partants} et les \omterm{def:b1:electronic-effects:intermediates}{carbocations} au
\cref{ch:b1:electronic-effects}~; les \omterm{def:b1:intermolecular-forces-solvents:solvent-classes}{solvants protiques} et aprotiques
au \cref{ch:b1:intermolecular-forces-solvents}.
\end{recall}

\section{La substitution et ses partenaires}

\begin{definition}[Substitution nucléophile, substrat]\label{def:b1:nucleophilic-substitution:substitution}
Une \emph{substitution nucléophile}\index{substitution nucleophile@substitution nucléophile} est le
remplacement, sur un carbone saturé, d'un \omterm{def:b1:electronic-effects:nucleophile}{groupe partant} Y par un
\omterm{def:b1:electronic-effects:nucleophile}{nucléophile} Nu~: $\mathrm{Nu^-} + \mathrm{R{-}Y} \to \mathrm{R{-}Nu} +
\mathrm{Y^-}$. Le composé R--Y est le \emph{substrat}\index{substrat}~;
le carbone qui porte Y est \omterm{def:b1:electronic-effects:nucleophile}{électrophile} parce que la liaison C--Y est
polarisée vers Y.
\end{definition}

Les halogénoalcanes sont les \omterm{def:b1:nucleophilic-substitution:substitution}{substrats} typiques, avec \ce{Cl-}, \ce{Br-}
ou \ce{I-} comme \omterm{def:b1:electronic-effects:nucleophile}{groupes partants}. L'hydroxyde donne des alcools~; les
alcoolates donnent des éthers~; le cyanure donne des nitriles (un carbone
de plus)~; l'ammoniac et les amines donnent des amines~; l'iodure
s'échange avec les autres halogénures. L'eau et les alcools, utilisés
comme solvants, peuvent aussi jouer le rôle de \omterm{def:b1:electronic-effects:nucleophile}{nucléophiles}~: la
substitution est alors une \omterm{def:b1:nucleophilic-substitution:sn1}{solvolyse}.

\section{Le mécanisme SN2}

Lorsque le bromométhane réagit avec l'hydroxyde dans l'eau, doubler
l'une ou l'autre concentration double la vitesse~:
$v = k[\ce{CH3Br}][\ce{OH-}]$, réaction d'ordre un par rapport à chaque
réactif.

\begin{definition}[Mécanisme SN2, inversion de Walden]\label{def:b1:nucleophilic-substitution:sn2}
Le \emph{mécanisme SN2}\index{mecanisme SN2@mécanisme SN2} (\omterm{def:b1:nucleophilic-substitution:substitution}{substitution nucléophile} bimoléculaire)
est un \omterm{def:b1:elementary-steps:elementary-step}{acte élémentaire} unique dans lequel le \omterm{def:b1:electronic-effects:nucleophile}{nucléophile} attaque le
carbone du côté opposé au \omterm{def:b1:electronic-effects:nucleophile}{groupe partant} pendant que celui-ci s'en va.
Les trois autres substituants du carbone basculent, comme un parapluie
qui se retourne~: c'est l'\emph{inversion de
Walden}\index{inversion de Walden}.
\end{definition}

\begin{omfigure}
\setchemfig{atom sep=2.1em}
\schemestart
\chemfig{H@{nu}O^{-}}
\arrow{0}[,0.35]
\chemfig{@{c}C(-[:150]H_3C)(<[:210]H)(<:[:250]C_2H_5)-[@{cb}]@{br}Br}
\arrow{->}
\chemfig{[:0]HO-C(-[:30]CH_3)(<[:-30]H)(<:[:-70]C_2H_5)}
\+
\chemfig{Br^{-}}
\schemestop
\chemmove{\draw[->, thick, shorten <=2pt, shorten >=3pt](nu) .. controls +(15:5mm) and +(175:5mm) .. (c);
          \draw[->, thick, shorten <=2pt](cb) .. controls +(90:6mm) and +(90:6mm) .. (br);}

\medskip
\begin{tikzpicture}[font=\small]
  \node[draw, dashed, rounded corners, inner sep=6pt] at (0,0)
    {$\left[\vcenter{\hbox{\ce{HO}$\cdots$\chemfig{C(-[2]CH_3)(-[:-30]H)(-[:210]C_2H_5)}$\cdots$\ce{Br}}}\right]^{-}$};
  \node[below] at (0,-1.1) {état de transition~: cinq groupes autour de C, dont trois dans un plan};
\end{tikzpicture}

{\small La réaction SN2 de l'hydroxyde avec le \cip{S}-2-bromobutane. Le
\omterm{def:b1:electronic-effects:nucleophile}{nucléophile} attaque la face opposée au brome (\omterm{def:b1:electronic-effects:curly-arrow}{flèches courbes})~; dans
l'\omterm{def:b1:elementary-steps:energy-profile}{état de transition}, la liaison O--C se forme pendant que la liaison
C--Br se rompt~; le produit est le \cip{R}-butan-2-ol~: les trois autres
groupes ont basculé de l'autre côté.}
\end{omfigure}

\begin{proposition}[Loi de vitesse de la SN2]\label{prop:b1:nucleophilic-substitution:sn2-law}
Pour une réaction SN2, $v = k[\mathrm{RY}][\mathrm{Nu}]$.
\end{proposition}

\begin{proof}
Le mécanisme est un \omterm{def:b1:elementary-steps:elementary-step}{acte élémentaire} bimoléculaire unique~; d'après la
loi d'action des masses appliquée aux \omterm{def:b1:elementary-steps:elementary-step}{actes élémentaires}
(\cref{ch:b1:elementary-steps}), sa vitesse est proportionnelle au
produit des concentrations des deux partenaires.
\end{proof}

\begin{proposition}[Stéréochimie de la SN2]\label{prop:b1:nucleophilic-substitution:sn2-inversion}
Une réaction SN2 sur un carbone stéréogène donne un seul stéréoisomère,
dans lequel la disposition des trois groupes inchangés est inversée~:
une réaction de ce type est stéréospécifique.
\end{proposition}

\begin{proof}
Le \omterm{def:b1:electronic-effects:nucleophile}{nucléophile} ne peut atteindre que le côté libre du carbone, opposé au
\omterm{def:b1:electronic-effects:nucleophile}{groupe partant}~; les trois autres groupes passent par un plan dans
l'\omterm{def:b1:elementary-steps:energy-profile}{état de transition} et finissent de l'autre côté. La nouvelle liaison
pointe là où l'ancienne ne pointait pas~: chaque molécule de \omterm{def:b1:nucleophilic-substitution:substitution}{substrat}
donne la disposition miroir de ses groupes. Le descripteur passe en
général de \cip{S} à \cip{R} ou l'inverse, mais pas toujours, car il
dépend des rangs du \omterm{def:b1:electronic-effects:nucleophile}{nucléophile} et du \omterm{def:b1:electronic-effects:nucleophile}{groupe partant}.
\end{proof}

\begin{definition}[Réaction stéréospécifique]\label{def:b1:nucleophilic-substitution:stereospecific}
Une réaction est une \emph{réaction stéréospécifique}\index{reaction stereospecifique@réaction stéréospécifique}
lorsque des \omterm{def:b1:nucleophilic-substitution:substitution}{substrats} \omterm{def:b1:stereochemistry-in-depth:isomers}{stéréoisomères} donnent des produits
\omterm{def:b1:stereochemistry-in-depth:isomers}{stéréoisomères} différents, chaque \omterm{def:b1:nucleophilic-substitution:substitution}{substrat} donnant son propre produit.
\end{definition}

\section{Le mécanisme SN1}

Le 2-bromo-2-méthylpropane réagit avec l'eau à une vitesse qui ne dépend
pas de la concentration du \omterm{def:b1:electronic-effects:nucleophile}{nucléophile}~: $v = k[(\ce{CH3})_3\ce{CBr}]$.
Une étape qui ne fait pas intervenir le \omterm{def:b1:electronic-effects:nucleophile}{nucléophile} doit fixer la
vitesse.

\begin{definition}[Mécanisme SN1, racémisation, solvolyse]\label{def:b1:nucleophilic-substitution:sn1}
Le \emph{mécanisme SN1}\index{mecanisme SN1@mécanisme SN1} (\omterm{def:b1:nucleophilic-substitution:substitution}{substitution nucléophile}
monomoléculaire) comporte deux étapes~: l'\omterm{def:b1:electronic-effects:cleavage}{hétérolyse} lente de la liaison
C--Y, qui donne un \omterm{def:b1:electronic-effects:intermediates}{carbocation}, puis l'addition rapide du \omterm{def:b1:electronic-effects:nucleophile}{nucléophile}
sur le \omterm{def:b1:electronic-effects:intermediates}{carbocation}. Une réaction qui transforme un énantiomère unique en
un mélange des deux est une \emph{racémisation}\index{racemisation@racémisation}~; une substitution
dans laquelle le solvant est le \omterm{def:b1:electronic-effects:nucleophile}{nucléophile} est une
\emph{solvolyse}\index{solvolyse}.
\end{definition}

\begin{omfigure}
\setchemfig{atom sep=2.1em}
\schemestart
\chemfig{C(-[2]CH_3)(-[4]H_3C)(-[6]CH_3)-[@{cb2}]@{br2}Br}
\arrow{->[lente]}
\chemfig{C^{+}(-[2]CH_3)(-[:210]H_3C)(-[:-30]CH_3)}
\+
\chemfig{Br^{-}}
\schemestop

\smallskip
\schemestart
\chemfig{C^{+}(-[2]CH_3)(-[:210]H_3C)(-[:-30]CH_3)}
\arrow{->[rapide][\ce{H2O}]}
\chemfig{C(-[2]CH_3)(-[4]H_3C)(-[6]CH_3)-OH_2^{+}}
\arrow{->[\ce{-H+}]}
\chemfig{C(-[2]CH_3)(-[4]H_3C)(-[6]CH_3)-OH}
\schemestop
\chemmove{\draw[->, thick, shorten <=2pt](cb2) .. controls +(90:6mm) and +(90:6mm) .. (br2);}

\medskip
\begin{tikzpicture}[font=\small]
  \fill[omProp!15, draw=omProp] (0,0.2) ellipse (0.2 and 0.85);
  \fill[omProp!15, draw=omProp] (0,-0.2) ellipse (0.2 and 0.85);
  \draw[thick] (0,0) -- (-1.1,-0.2) node[left] {\ce{CH3}};
  \draw[thick] (0,0) -- (0.95,-0.45) node[right] {\ce{C2H5}};
  \draw[thick] (0,0) -- (0.4,0.5) node[above right] {H};
  \node[fill=white, inner sep=1pt] at (0,0) {\ce{C+}};
  \draw[->, thick, omDef] (-0.9,1.6) node[left] {\ce{H2O}} -- (-0.15,0.95);
  \draw[->, thick, omDef] (0.9,-1.6) node[right] {\ce{H2O}} -- (0.15,-0.95);
  \node[align=center] at (4.6,0) {attaque sur l'une ou l'autre face~:\\ les deux configurations,\\ racémique si rien d'autre\\ ne distingue les faces};
\end{tikzpicture}

{\small \omterm{def:b1:nucleophilic-substitution:sn1}{Solvolyse} SN1 du 2-bromo-2-méthylpropane dans l'eau~: ionisation
lente en \omterm{def:b1:electronic-effects:intermediates}{carbocation} tertiaire, puis capture par l'eau et perte d'un
proton. En bas~: un \omterm{def:b1:electronic-effects:intermediates}{carbocation} secondaire \omterm{def:b1:stereochemistry-in-depth:stereocentre}{chiral} (issu du
2-bromobutane) est plan, et l'eau s'additionne sur l'une ou l'autre face
avec la même probabilité.}
\end{omfigure}

\begin{proposition}[Loi de vitesse de la SN1]\label{prop:b1:nucleophilic-substitution:sn1-law}
Pour le \omterm{def:b1:nucleophilic-substitution:sn1}{mécanisme SN1}, si le \omterm{def:b1:electronic-effects:intermediates}{carbocation} est capturé beaucoup plus vite
qu'il ne revient au \omterm{def:b1:nucleophilic-substitution:substitution}{substrat}, $v = k_1[\mathrm{RY}]$, indépendamment du
\omterm{def:b1:electronic-effects:nucleophile}{nucléophile}.
\end{proposition}

\begin{proof}
Étapes~: $\mathrm{RY} \rightleftharpoons \mathrm{R^+} + \mathrm{Y^-}$
($k_1$,
$k_{-1}$), $\mathrm{R^+} + \mathrm{Nu} \to \mathrm{RNu}$ ($k_2$). Le
\omterm{def:b1:electronic-effects:intermediates}{carbocation} est un \omterm{def:b1:electronic-effects:intermediates}{intermédiaire réactif}~: d'après l'\omterm{def:b1:elementary-steps:qssa}{approximation des
états quasi stationnaires} (\cref{ch:b1:elementary-steps}), $k_1[\mathrm{RY}] =
k_{-1}[\mathrm{R^+}][\mathrm{Y^-}] + k_2[\mathrm{R^+}][\mathrm{Nu}]$, donc
$v = k_2[\mathrm{R^+}][\mathrm{Nu}] = \dfrac{k_1k_2[\mathrm{RY}][\mathrm{Nu}]}
{k_{-1}[\mathrm{Y^-}] + k_2[\mathrm{Nu}]}$. Lorsque $k_2[\mathrm{Nu}] \gg
k_{-1}[\mathrm{Y^-}]$, cette expression se réduit à $k_1[\mathrm{RY}]$~:
la première étape est cinétiquement déterminante.
\end{proof}

\begin{proposition}[Stéréochimie de la SN1]\label{prop:b1:nucleophilic-substitution:sn1-stereo}
Une réaction SN1 sur un carbone stéréogène donne les deux
\omterm{def:b1:stereochemistry-in-depth:isomers}{configurations}~: la réaction n'est pas stéréospécifique, et conduit à
une \omterm{def:b1:nucleophilic-substitution:sn1}{racémisation} lorsque rien ne distingue les deux faces du
\omterm{def:b1:electronic-effects:intermediates}{carbocation}.
\end{proposition}

\begin{proof}
Le \omterm{def:b1:electronic-effects:intermediates}{carbocation} est plan (\cref{ch:b1:electronic-effects})~: ses deux
faces sont images l'une de l'autre. Si le \omterm{def:b1:electronic-effects:nucleophile}{groupe partant} s'est
suffisamment éloigné avant la capture, le \omterm{def:b1:electronic-effects:nucleophile}{nucléophile} s'additionne sur
l'une ou l'autre face à la même vitesse, ce qui donne les deux
\omterm{def:b1:stereochemistry-in-depth:enantiomers}{énantiomères} en quantités égales. En pratique, l'ion partant masque
encore une face pendant un moment, et l'on observe souvent un léger
excès d'inversion.
\end{proof}

\begin{omfigure}
\begin{tikzpicture}
\begin{axis}[
  width=0.8\linewidth, height=5.0cm, xmin=0, xmax=10, ymin=-30, ymax=110,
  axis x line=bottom, axis y line=left, xlabel={coordonnée réactionnelle},
  ylabel={énergie}, xtick=\empty, ytick=\empty,
  label style={font=\small}, legend style={font=\scriptsize, draw=none, at={(0.98,0.98)}, anchor=north east}]
\addplot[omDef, very thick] table[x=x, y=E] {figdata/out/bachelor-1-energy-profiles-sn2.dat};
\addplot[omProp, very thick, dashed] table[x=x, y=E] {figdata/out/bachelor-1-energy-profiles-sn1.dat};
\legend{SN2~: une étape, SN1~: deux étapes}
\node[font=\scriptsize, anchor=north] at (axis cs:5,68) {carbocation};
\node[font=\scriptsize, anchor=north] at (axis cs:0.6,-3) {RY};
\node[font=\scriptsize, anchor=south] at (axis cs:9.6,-18) {RNu};
\end{axis}
\end{tikzpicture}

{\small \omterm{def:b1:elementary-steps:energy-profile}{Profils énergétiques} schématiques. SN2~: un seul \omterm{def:b1:elementary-steps:energy-profile}{état de
transition}, cinq groupes autour du carbone. SN1~: une première barrière,
plus haute (l'ionisation, \omterm{def:b1:elementary-steps:rds}{étape cinétiquement déterminante}), le
\omterm{def:b1:electronic-effects:intermediates}{carbocation} dans un puits peu profond, puis une petite barrière vers sa
capture.}
\end{omfigure}

\section{Ce qui détermine le mécanisme}

\begin{proposition}[Le substrat]\label{prop:b1:nucleophilic-substitution:substrate}
Les \omterm{def:b1:nucleophilic-substitution:substitution}{substrats} méthyliques et primaires réagissent par SN2~; les
\omterm{def:b1:nucleophilic-substitution:substitution}{substrats} tertiaires par SN1~; les \omterm{def:b1:nucleophilic-substitution:substitution}{substrats} secondaires par l'un ou
l'autre, selon les autres facteurs. Les \omterm{def:b1:nucleophilic-substitution:substitution}{substrats} allyliques et
benzyliques réagissent rapidement par les deux.
\end{proposition}

\begin{proof}
La SN2 exige de la place derrière le carbone~: chaque groupe alkyle qui
remplace un hydrogène encombre l'\omterm{def:b1:elementary-steps:energy-profile}{état de transition}, qui porte cinq
groupes, et ralentit la réaction~; trois groupes la bloquent. La SN1
exige un \omterm{def:b1:electronic-effects:intermediates}{carbocation} stable~: l'ordre de stabilité
(\cref{prop:b1:electronic-effects:carbocation-order}) rend l'ionisation
rapide pour les \omterm{def:b1:electronic-effects:intermediates}{carbocations} tertiaires, possible pour les secondaires,
et trop lente pour les primaires et le méthyle. Les \omterm{def:b1:electronic-effects:intermediates}{carbocations}
allyliques et benzyliques sont stabilisés par mésomérie, et le même
système $\pi$ stabilise l'\omterm{def:b1:elementary-steps:energy-profile}{état de transition} de la SN2.
\end{proof}

\begin{proposition}[Le groupe partant]\label{prop:b1:nucleophilic-substitution:leaving-group}
Les deux mécanismes sont d'autant plus rapides que le \omterm{def:b1:electronic-effects:nucleophile}{groupe partant} est
meilleur, c'est-à-dire que la base qu'il devient est plus faible~:
$\ce{I-} > \ce{Br-} > \ce{Cl-} \gg \ce{F-}$~; \ce{OH-} et \ce{RO-} ne
partent pas.
\end{proposition}

\begin{proof}
Dans les deux mécanismes, la liaison C--Y est rompue au cours de l'\omterm{def:b1:elementary-steps:rds}{étape
cinétiquement déterminante} ou avant elle, et l'\omterm{def:b1:elementary-steps:energy-profile}{état de transition} porte
une partie de la charge négative sur Y. Un groupe qui porte bien une
charge négative --- la base conjuguée d'un \omterm{def:b1:predominant-reaction:strong-acid}{acide fort} --- abaisse cet
\omterm{def:b1:elementary-steps:energy-profile}{état de transition}. \ce{HI}, \ce{HBr} et \ce{HCl} sont des \omterm{def:b1:predominant-reaction:strong-acid}{acides forts},
\ce{HF} un \omterm{def:b1:predominant-reaction:strong-acid}{acide faible} ($\mathrm{p}K_a$ 3.2), l'eau et les alcools des
acides très faibles~: \ce{OH-} et \ce{RO-} sont des \omterm{def:b1:predominant-reaction:strong-acid}{bases fortes} et de
mauvais \omterm{def:b1:electronic-effects:nucleophile}{groupes partants}.
\end{proof}

\begin{proposition}[Le solvant]\label{prop:b1:nucleophilic-substitution:solvent}
Les \omterm{def:b1:intermolecular-forces-solvents:solvent-classes}{solvants polaires} protiques (eau, alcools) favorisent la SN1~; les
\omterm{def:b1:intermolecular-forces-solvents:solvent-classes}{solvants polaires} aprotiques (propanone, diméthylsulfoxyde,
diméthylformamide) favorisent la SN2 avec les \omterm{def:b1:electronic-effects:nucleophile}{nucléophiles} anioniques.
\end{proposition}

\begin{proof}
La SN1 crée deux ions à partir d'un \omterm{def:b1:nucleophilic-substitution:substitution}{substrat} neutre~: un \omterm{def:b1:intermolecular-forces-solvents:solvent-classes}{solvant polaire}
protique les stabilise tous deux, le cation par ses \omterm{def:b1:lewis-resonance-vsepr:lewis}{doublets non liants}
et l'anion par \omterm{def:b1:intermolecular-forces-solvents:hbond}{liaisons hydrogène}, et abaisse la barrière d'ionisation.
Dans la SN2, la charge d'un \omterm{def:b1:electronic-effects:nucleophile}{nucléophile} anionique est étalée dans l'\omterm{def:b1:elementary-steps:energy-profile}{état
de transition}~; un \omterm{def:b1:intermolecular-forces-solvents:solvent-classes}{solvant protique}, qui lie fortement le petit anion
par \omterm{def:b1:intermolecular-forces-solvents:hbond}{liaisons hydrogène}, la ralentit, alors qu'un \omterm{def:b1:intermolecular-forces-solvents:solvent-classes}{solvant aprotique}
laisse l'anion libre et réactif
(\cref{ch:b1:intermolecular-forces-solvents}).
\end{proof}

\begin{method}[Choisir le mécanisme]\label{met:b1:nucleophilic-substitution:choose-mechanism}
\begin{enumerate}
  \item \omterm{def:b1:nucleophilic-substitution:substitution}{Substrat}~: méthylique ou primaire $\to$ SN2~; tertiaire $\to$
        SN1~; secondaire $\to$ poursuivre.
  \item \omterm{def:b1:electronic-effects:nucleophile}{Nucléophile}~: fort et anionique ($\ce{OH-}$, $\ce{RO-}$,
        $\ce{CN-}$,
        $\ce{I-}$) $\to$ SN2~; faible et neutre (eau, alcool) $\to$ SN1.
  \item Solvant~: aprotique $\to$ SN2~; protique $\to$ SN1.
  \item Prévoir la stéréochimie (inversion ou \omterm{def:b1:nucleophilic-substitution:sn1}{racémisation}) et vérifier
        qu'une \omterm{def:b1:predominant-reaction:strong-acid}{base forte} ne provoque pas plutôt une élimination
        (\cref{ch:b1:elimination}).
\end{enumerate}
\end{method}

\begin{omfigure}
\begin{tabular}{lll}
\hline
 & SN2 & SN1 \\
\hline
étapes & une & deux, via un \omterm{def:b1:electronic-effects:intermediates}{carbocation} \\
\omterm{def:b1:rate-laws:order}{loi de vitesse} & $k[\mathrm{RY}][\mathrm{Nu}]$ & $k[\mathrm{RY}]$ \\
\omterm{def:b1:nucleophilic-substitution:substitution}{substrat} & méthyle $>$ primaire $>$ secondaire & tertiaire $>$ secondaire \\
\omterm{def:b1:electronic-effects:nucleophile}{nucléophile} & fort, anionique & faible, neutre (souvent le solvant) \\
solvant & polaire aprotique & polaire protique \\
stéréochimie & inversion (stéréospécifique) & \omterm{def:b1:nucleophilic-substitution:sn1}{racémisation} (en général) \\
\hline
\end{tabular}

\smallskip
{\small Les deux mécanismes côte à côte.}
\end{omfigure}

\section{Exercices}

\begin{exercise}[$\star$]\label{exo:b1:nucleophilic-substitution:1}
Écrire les produits~: du bromoéthane avec l'hydroxyde de sodium~; de
l'iodométhane avec l'éthanolate de sodium~; du 1-chloropropane avec le
cyanure de potassium~; du bromométhane avec l'ammoniac (première étape).
Identifier le \omterm{def:b1:nucleophilic-substitution:substitution}{substrat}, le \omterm{def:b1:electronic-effects:nucleophile}{nucléophile} et le \omterm{def:b1:electronic-effects:nucleophile}{groupe partant}.
\end{exercise}

\begin{exercise}[$\star$]\label{exo:b1:nucleophilic-substitution:2}
Pour une réaction $\mathrm{RY} + \mathrm{Nu}$ de vitesse
$v = k[\mathrm{RY}][\mathrm{Nu}]$, que devient la \omterm{def:b1:rate-laws:initial-rate}{vitesse initiale} si
l'on triple $[\mathrm{Nu}]$~? Si l'on divise par deux les deux
concentrations~? Mêmes questions si $v = k[\mathrm{RY}]$.
\end{exercise}

\begin{exercise}[$\star$]\label{exo:b1:nucleophilic-substitution:3}
Prévoir le mécanisme (SN1 ou SN2)~: 1-bromobutane avec l'iodure de
sodium dans la propanone~; 2-bromo-2-méthylpropane dans l'eau~;
bromométhane avec l'hydroxyde~; 2-bromobutane dans l'éthanol sans base
ajoutée.
\end{exercise}

\begin{exercise}[$\star$]\label{exo:b1:nucleophilic-substitution:4}
Dessiner en \omterm{def:b1:stereochemistry-in-depth:representations}{représentation de Cram} le produit de la réaction SN2 du
cyanure sur le \cip{R}-2-bromobutane et donner son descripteur.
\end{exercise}

\begin{exercise}[$\star\star$]\label{exo:b1:nucleophilic-substitution:5}
Écrire le \omterm{def:b1:nucleophilic-substitution:sn1}{mécanisme SN1} complet de l'hydrolyse du
2-bromo-2-méthylpropane, avec les \omterm{def:b1:electronic-effects:curly-arrow}{flèches courbes}, y compris la perte du
proton. Pourquoi la \omterm{def:b1:rate-laws:order}{loi de vitesse} est-elle d'ordre un alors que trois
espèces interviennent~?
\end{exercise}

\begin{exercise}[$\star\star$]\label{exo:b1:nucleophilic-substitution:6}
Expliquer pourquoi le 1-bromo-2,2-diméthylpropane, un bromure primaire,
réagit très lentement par SN2 et ne réagit pas par SN1.
\end{exercise}

\begin{exercise}[$\star\star$]\label{exo:b1:nucleophilic-substitution:7}
Le \cip{S}-2-iodooctane perd son \omterm{def:b1:stereochemistry-in-depth:optical-activity}{activité optique} lorsqu'on le conserve
dans la propanone avec de l'iodure de sodium, bien qu'aucun composé
nouveau ne se forme. Expliquer, et dire pourquoi la vitesse de perte de
l'\omterm{def:b1:stereochemistry-in-depth:optical-activity}{activité optique} est le double de la vitesse de substitution.
\end{exercise}

\begin{exercise}[$\star\star$]\label{exo:b1:nucleophilic-substitution:8}
Classer par vitesse de SN2 avec l'iodure~: 2-bromo-2-méthylpropane,
bromométhane, 2-bromopropane, 1-bromopropane. Classer les mêmes composés
par vitesse de \omterm{def:b1:nucleophilic-substitution:sn1}{solvolyse} SN1.
\end{exercise}

\begin{exercise}[$\star\star$]\label{exo:b1:nucleophilic-substitution:9}
Un bromure secondaire optiquement pur subit une \omterm{def:b1:nucleophilic-substitution:sn1}{solvolyse}~; le produit
fait tourner la lumière de \qty{12}{\percent} de la valeur attendue pour
une inversion pure. Donner les proportions des deux \omterm{def:b1:stereochemistry-in-depth:enantiomers}{énantiomères} du
produit et interpréter.
\end{exercise}

\begin{exercise}[$\star\star\star$]\label{exo:b1:nucleophilic-substitution:10}
Établir la \omterm{def:b1:rate-laws:order}{loi de vitesse} générale de la SN1 par l'\omterm{def:b1:elementary-steps:qssa}{approximation des
états quasi stationnaires}, et montrer qu'ajouter dès le départ l'ion
partant \ce{Y-} ralentit la réaction (l'\omterm{def:b1:precipitation:common-ion}{effet d'ion commun} de la
cinétique). À quelle condition cet effet est-il négligeable~?
\end{exercise}

\begin{exercise}[$\star\star\star$]\label{exo:b1:nucleophilic-substitution:11}
Deux \omterm{def:b1:rate-laws:half-life}{temps de demi-réaction} pour la réaction du bromométhane avec
l'hydroxyde~: avec $[\ce{OH-}]_0 = [\ce{CH3Br}]_0 = C_0$, montrer que
$1/[\ce{CH3Br}] = 1/C_0 +
kt$ et donner $t_{1/2}$. Avec $[\ce{OH-}]_0 = 50\,C_0$, montrer que la
décroissance est exponentielle et donner $t_{1/2}$. Données de
l'exercice~: $k =
\qty{2.0e-4}{L/(mol.s)}$, $C_0 = \qty{0.010}{mol/L}$.
\end{exercise}

\begin{exercise}[$\star\star\star$]\label{exo:b1:nucleophilic-substitution:12}
Expliquer pourquoi un alcool, \ce{ROH}, ne réagit pas avec le bromure de
sodium, mais réagit avec l'acide bromhydrique concentré pour donner
\ce{RBr}. Écrire le mécanisme pour un alcool tertiaire et pour un alcool
primaire.
\end{exercise}

\section{Problème~: la solvolyse du chlorure de \textit{tert}-butyle}

\begin{problem}[{Problème du week-end --- ordre et constante de vitesse
à partir de mesures de conductivité, temps de demi-réaction, mécanisme
et stéréochimie, et énergie d'activation de la solvolyse}]
\label{pb:b1:nucleophilic-substitution:1}
On dissout du 2-chloro-2-méthylpropane (chlorure de \textit{tert}-butyle),
à faible concentration, dans un mélange eau--propanone. Sa \omterm{def:b1:nucleophilic-substitution:sn1}{solvolyse}
\ce{(CH3)3CCl + H2O -> (CH3)3COH + H+ + Cl-} produit des ions, et l'on
enregistre la conductivité $\kappa$ de la solution, proportionnelle à la
quantité d'ions formés. Au bout d'un temps très long, elle atteint
$\kappa_\infty = \qty{2.000}{mS/cm}$ aux deux températures. Mesures
(\unit{mS/cm})~:

\begin{center}
\begin{tabular}{lcccccccc}
\hline
$t$ (min) & 0 & 5 & 10 & 15 & 20 & 30 & 45 & 60 \\
\hline
$\kappa$, \qty{25.0}{\celsius} & 0 & 0.172 & 0.329 & 0.473 & 0.605 & 0.835 & 1.110 & 1.321 \\
$\kappa$, \qty{35.0}{\celsius} & 0 & 0.507 & 0.885 & 1.168 & 1.379 & 1.654 & 1.856 & 1.940 \\
\hline
\end{tabular}
\end{center}
% data generated in figdata/bachelor-1/solvolysis.py (story data)

\medskip
\noindent\textbf{Partie I --- L'ordre.}
\begin{enumerate}
  \item Pourquoi la conductivité est-elle proportionnelle à la quantité
        de \omterm{def:b1:nucleophilic-substitution:substitution}{substrat} qui a réagi~?
  \item Exprimer $[\mathrm{RCl}]/[\mathrm{RCl}]_0$ en fonction de $\kappa$
        et de $\kappa_\infty$.
  \item Écrire la loi intégrée d'une réaction d'ordre 1 et la grandeur à
        tracer en fonction de $t$ pour la tester.
  \item Calculer $\ln(\kappa_\infty - \kappa)$ à \qty{25.0}{\celsius}
        pour chaque instant.
  \item Le tracé est-il une droite~? Conclure sur l'ordre.
  \item L'eau est en large excès. Qu'est-ce que cela masque dans
        l'ordre~?
\end{enumerate}

\medskip
\noindent\textbf{Partie II --- \omterm{def:b1:rate-laws:order}{Constante de vitesse} et \omterm{def:b1:rate-laws:half-life}{temps de demi-réaction}.}
\begin{enumerate}[resume]
  \item Déterminer la \omterm{def:b1:rate-laws:order}{constante de vitesse} à \qty{25.0}{\celsius}, en
        \unit{s^{-1}}.
  \item Calculer le \omterm{def:b1:rate-laws:half-life}{temps de demi-réaction}.
  \item À quel instant \qty{90}{\percent} du \omterm{def:b1:nucleophilic-substitution:substitution}{substrat} est-il consommé~?
  \item Déterminer la \omterm{def:b1:rate-laws:order}{constante de vitesse} à \qty{35.0}{\celsius}.
  \item De quel facteur la vitesse a-t-elle augmenté pour dix degrés~?
  \item Vérifier une valeur de la série à \qty{35}{\celsius} à l'aide de
        la loi.
\end{enumerate}

\medskip
\noindent\textbf{Partie III --- Le mécanisme.}
\begin{enumerate}[resume]
  \item Quel mécanisme le \omterm{def:b1:nucleophilic-substitution:substitution}{substrat} et le solvant suggèrent-ils~?
  \item L'écrire avec des \omterm{def:b1:electronic-effects:curly-arrow}{flèches courbes}.
  \item Montrer qu'il est en accord avec l'ordre trouvé.
  \item L'ajout d'hydroxyde de sodium à \qty{0.01}{mol/L} change à peine
        la \omterm{def:b1:rate-laws:rate}{vitesse de disparition} du chlorure. Pourquoi~?
  \item La même expérience menée avec le chlorure tertiaire \omterm{def:b1:stereochemistry-in-depth:stereocentre}{chiral}
        3-chloro-3-méthylhexane, optiquement pur, donne un alcool
        d'\omterm{def:b1:stereochemistry-in-depth:optical-activity}{activité optique} presque nulle. Expliquer.
  \item Tracer le \omterm{def:b1:elementary-steps:energy-profile}{profil énergétique} de la réaction et y repérer l'\omterm{def:b1:elementary-steps:rds}{étape
        cinétiquement déterminante}.
  \item Pourquoi le 1-chlorobutane réagirait-il à peine dans les mêmes
        conditions~?
\end{enumerate}

\medskip
\noindent\textbf{Partie IV --- L'\omterm{def:b1:rate-laws:activation-energy}{énergie d'activation}.}
\begin{enumerate}[resume]
  \item Écrire la \omterm{thm:b1:rate-laws:arrhenius}{loi d'Arrhenius}.
  \item Exprimer $E_a$ à partir des deux \omterm{def:b1:rate-laws:order}{constantes de vitesse}.
  \item Calculer $E_a$ ($R = \qty{8.314}{J/(K.mol)}$).
  \item Que mesure cette énergie, dans le mécanisme~?
  \item Prévoir la \omterm{def:b1:rate-laws:order}{constante de vitesse} à \qty{45.0}{\celsius}.
  \item Donner l'\omterm{def:b1:rate-laws:activation-energy}{énergie d'activation} de la \omterm{def:b1:nucleophilic-substitution:sn1}{solvolyse}, à deux chiffres
        significatifs.
\end{enumerate}
\end{problem}
